
In recent years, rapid urbanization and the increasing number of vehicles have made parking management a major challenge in residential, commercial, institutional, and public areas. Traditional parking systems usually depend on manual gate control, manual slot checking, and manual record keeping. This process is time-consuming, less accurate, and inefficient when the number of vehicles is high.

Automation, robotics, embedded systems, and Internet of Things (IoT) technologies can improve parking management by detecting vehicles, verifying access, controlling gates, updating available slots, and storing parking data automatically. In this project, an \textbf{Automated Car Parking System} is designed using Arduino UNO, IR sensors, RFID, LCD display, ESP-01S Wi-Fi communication, and a servo motor-controlled two-link robotic arm gate mechanism.

The system detects a vehicle at the entrance using an IR sensor and checks parking slot availability. If a slot is available, the user scans an RFID card for authentication. After successful verification, the Arduino controls the servo motor to move the two-link robotic arm gate. The available slot count is updated on the LCD and the entry record is sent to an online server. During exit, the system detects the vehicle, opens the robotic arm gate, increases the available slot count, and stores the exit record. Thus, the proposed system combines robotic gate control and IoT-based monitoring to create a low-cost smart parking prototype.

\section{Background and Motivation}

Parking space management has become difficult in busy places such as shopping malls, offices, universities, hospitals, and residential buildings. Drivers often waste time searching for available parking spaces, which increases congestion near parking areas. Manual systems may also create errors in slot counting and vehicle entry-exit records.

The motivation of this project is to develop a simple, low-cost, and automated parking system suitable for academic prototype implementation. The system uses IR sensors for vehicle detection, RFID for authorized access, an LCD for real-time slot display, ESP-01S for online data transmission, and a two-link robotic arm as the gate mechanism. The robotic arm adds a robotics-based control feature and demonstrates practical use of actuator movement in parking automation.

\section{Problem Statement}

Traditional parking systems have several limitations. Manual gate operation can cause delays during vehicle entry and exit. Manual slot counting can be inaccurate, and drivers may not know whether parking space is available before reaching the gate. Manual record keeping is also difficult to maintain for a long time.

Security is another concern in controlled parking areas. Without authentication, unauthorized vehicles may enter the parking area. Therefore, an automated system is needed that can detect vehicles, verify authorized users, control the gate, update slot availability, and store parking data using IoT communication.

\section{Objectives of the Project}

The main objective of this project is to design and implement an Automated Car Parking System that can control vehicle entry and exit using sensors, RFID, microcontroller-based control, IoT communication, and a two-link robotic arm gate.

The specific objectives are:

\begin{enumerate}
    \item To design a prototype automated parking system using Arduino UNO.
    \item To detect vehicles at entry and exit points using IR sensors.
    \item To use RFID for authorized vehicle entry.
    \item To control a servo motor-based two-link robotic arm gate.
    \item To display available parking slots on a 16x2 I2C LCD.
    \item To update parking slot count automatically after entry and exit.
    \item To send entry and exit records to an online server using ESP-01S.
    \item To demonstrate the integration of robotics and IoT in parking automation.
\end{enumerate}

\section{Research Questions}

The project is guided by the following research questions:

\begin{enumerate}
    \item How can an Arduino-based system automate vehicle entry and exit?
    \item How effectively can IR sensors detect vehicle presence?
    \item How can RFID be used for authorized parking access?
    \item Can a servo-controlled two-link robotic arm operate as a reliable gate?
    \item How can parking slot information be updated and stored automatically?
\end{enumerate}

\section{Scope of the Work}

This project focuses on a prototype-level Automated Car Parking System. The system uses Arduino UNO, IR sensors, RFID module, ESP-01S Wi-Fi module, LCD display, and a servo motor-controlled two-link robotic arm. The main functions include vehicle detection, RFID authentication, robotic gate control, parking slot counting, LCD-based status display, and IoT-based data transmission.

The scope is limited to a small-scale lab prototype. It does not include automatic number plate recognition, camera-based detection, payment system, mobile reservation, GPS guidance, or large-scale commercial deployment. However, these features can be added in future versions.

\section{Expected Outcomes}

The expected outcome is a functional prototype that can detect vehicles, verify RFID cards, operate the two-link robotic arm gate, update slot availability, display parking status, and send parking records online. The system is expected to reduce manual effort, improve parking efficiency, and demonstrate the practical application of robotics and IoT in smart parking management.

\section{Impacts of the Project}

The proposed system can improve parking management by reducing manual gate operation and manual slot counting. It can save time for drivers, reduce congestion near parking gates, and support more organized parking control. Since the system stores parking data online, data privacy and secure communication should be considered in real deployment.

\subsection{Impact on Societal and Cultural Issues}

The system can support society by making parking more efficient and reducing waiting time. It also encourages the use of automation and IoT-based solutions in daily life, supporting smart city development and technology awareness.

\subsection{Impact on Health, Safety, and Legal Issues}

The automated gate operation can reduce human error and improve safety during vehicle entry and exit. The system has no direct health impact, but it can reduce driver stress by making parking faster. From a legal perspective, RFID and entry-exit data should be protected from unauthorized access.

\subsection{Impact on Environment and Sustainability Issues}

Efficient parking management can reduce unnecessary vehicle movement and waiting time, which may help lower fuel consumption and emissions. The system also uses low-cost and low-power electronic components suitable for sustainable prototype development.

\section{Report Organization}

The report is organized as follows:

\textbf{Chapter 1: Introduction}\\
This chapter presents the background, motivation, problem statement, objectives, scope, expected outcomes, impacts, and report organization.

\textbf{Chapter 2: Related Works}\\
This chapter reviews existing smart parking systems, IoT-based parking management, RFID access control, and automated gate systems.

\textbf{Chapter 3: System Analysis and Design}\\
This chapter describes requirements, system architecture, hardware setup, circuit connection, and robotic arm gate structure.

\textbf{Chapter 4: Methodology and Implementation}\\
This chapter explains the development method, tools, data handling, algorithm, flowchart, and module implementation.

\textbf{Chapter 5: Results and Evaluation}\\
This chapter presents testing setup, performance metrics, result analysis, discussion, and limitations.

\textbf{Chapter 6: Project Planning and Budget}\\
This chapter includes the project timeline, cost estimation, and resource planning.

\textbf{Chapter 7: Standards and Design Constraints}\\
This chapter discusses standards, design constraints, CEP, CEA, and knowledge profile mapping.

\textbf{Chapter 8: Conclusion and Future Work}\\
This chapter summarizes the project achievements, limitations, and future improvement possibilities.

